% Compile with pdfLaTeX. Required packages: amsmath, amssymb, algorithm,
% algpseudocode, and geometry.
\documentclass[11pt]{article}
\usepackage{amsmath,amssymb}
\usepackage{algorithm}
\usepackage{algpseudocode}
\usepackage[a4paper,margin=1in]{geometry}

\title{Selective Frequency-Domain Backdoor Mitigation}
\author{}
\date{}

\begin{document}
\maketitle

\section*{Key notation}
\noindent\fbox{\begin{minipage}{0.94\linewidth}
\small
\begin{itemize}
    \item $\mathcal{D}$: clean training set; $x\in[0,1]^{C\times H\times W}$ is an image.
    \item $\rho$: poisoning ratio; $y^*$: attacker-selected target class.
    \item $\alpha$: trigger strength; $f_s$: suspicious classifier.
    \item $G_\theta$: spectral-correction generator; $f_r$: repaired classifier.
    \item $A$: FFT amplitude, $P$: FFT phase, and $M$: learned correction gate.
    \item $A_c$: clean amplitude, $A_t$: triggered amplitude, and
    $A_{corr}$: corrected amplitude.
\end{itemize}
\end{minipage}}

\section*{Trigger construction}
For the sinusoidal experiment, the spatial trigger is
\begin{equation}
    S(u,v)=\sin\left(2\pi\left(\frac{f_xu}{W}+\frac{f_yv}{H}\right)\right),
    \qquad
    t=\operatorname{clip}(x+\alpha S,0,1).
\end{equation}

\section*{How the trigger strength $\alpha$ is selected}
There is no universal theorem stating that one value of $\alpha$ is correct for
every dataset, image size, classifier, or trigger type. The value is a
hyperparameter that controls the security--utility trade-off. A larger value
usually makes the attack easier for the classifier to learn, but it can also
make the modification more visible and can reduce clean-image performance.

For the sinusoidal trigger, the signal is bounded by $-1\leq S\leq 1$.
Ignoring pixel clipping, this gives
\begin{equation}
    \|t-x\|_\infty\leq\alpha,
    \qquad
    \mathbb{E}\left[(t-x)^2\right]\approx\frac{\alpha^2}{2}
\end{equation}
for a zero-mean sinusoid. Thus, $\alpha$ has a direct mathematical meaning:
it bounds the maximum added signal and its energy grows approximately with
$\alpha^2$. Clipping and the content of the original image make the measured
pixel distortion image-dependent, so these expressions are design guides and
not substitutes for reporting measured PSNR, SSIM, or pixel error.

The values used in the experiments form an exploratory low-to-high grid:
\begin{center}
\small
\begin{tabular}{c|c|l}
\textbf{Experiment} & \textbf{Values of $\alpha$} & \textbf{Purpose}\\
\hline
CIFAR-100 baseline/ablation & $0.04,\ 0.08,\ 0.12$ & weak, moderate, and stronger trigger\\
STL-10 strength study & $0.03,\ 0.15,\ 0.25$ & broader resolution-specific stress test\\
FIBA amplitude study & $0.15,\ 0.30,\ 0.50$ & low, medium, and strong amplitude blending\\
\end{tabular}
\end{center}

The CIFAR-100 value $\alpha=0.08$ was used as the main baseline because it is
large enough to produce a measurable backdoor while remaining a controlled
perturbation. The neighbouring values $0.04$ and $0.12$ test whether the
defense is sensitive to a weaker or stronger sinusoidal signal. In the CIFAR
ablation, $0.08$ produced the lowest generator-corrected ASR among the tested
values while retaining a high suspicious ASR, so it was a reasonable
representative operating point rather than a claimed universal default.

FIBA uses a different scale. Inside the frequency mask $W$,
\begin{equation}
    A_t-A_c=\alpha W\odot(A_q-A_c).
\end{equation}
Consequently, $\alpha$ is an amplitude blending coefficient, not a direct
pixel-wise perturbation bound. The FIBA values $0.15$, $0.30$, and $0.50$ test
increasing proportions of the reference amplitude. They must not be compared
numerically with the sinusoidal values. The mask radius is a separate variable
that controls how large a frequency region is modified.

For a rigorous study, $\alpha$ should be selected on a validation set using a
predefined rule such as
\begin{equation}
    \alpha^*=\min_{\alpha}\left\{\alpha:\operatorname{ASR}(\alpha)\geq\tau,
    \operatorname{CleanAcc}(\alpha)\geq\gamma,
    \operatorname{Distortion}(\alpha)\leq\delta\right\}.
\end{equation}
The test set should then be used only once for final reporting. This makes the
choice reproducible and prevents selecting a trigger strength after seeing the
final test result.

\begin{algorithm}[t]
\caption{Selective spectral-correction defense}
\label{alg:selective-defense}
\begin{algorithmic}[1]
\Require Clean set $\mathcal{D}$; poisoning ratio $\rho$; target class $y^*$;
trigger settings; learning rates; loss weights.
\Ensure $f_s$, $G_\theta$, $f_r$, and evaluation metrics.
\Statex \textbf{One-line summary:} Poison data $\rightarrow$ train suspicious
model $\rightarrow$ analyze FFT spectra $\rightarrow$ predict gate $M$
$\rightarrow$ correct amplitude $\rightarrow$ reconstruct image
$\rightarrow$ repair model $\rightarrow$ evaluate.

\State Select approximately $\rho|\mathcal{D}|$ non-target training images.
\For{each selected image $(x_i,y_i)$}
    \State $t_i\gets\operatorname{Trigger}(x_i;\alpha)$.
    \State Add $(t_i,y^*)$ to the poisoned training set.
\EndFor
\State Train $f_s$ on clean samples and poisoned samples.
\State Measure clean accuracy and triggered attack success rate of $f_s$.

\For{each generator training epoch}
    \For{each clean minibatch $(x,y)$ with $y\neq y^*$}
        \State $t\gets\operatorname{Trigger}(x;\alpha)$.
        \Statex \textit{Compare the clean and triggered images in the frequency domain.}
        \State $(A_c,P_c)\gets\operatorname{FFT}(x)$ and
        $(A_t,P_t)\gets\operatorname{FFT}(t)$.
        \Statex \textit{The generator receives clean amplitude, triggered amplitude, and their difference.}
        \State $Z\gets\operatorname{Concat}(\mathcal{N}(\log(1+A_c)),
        \mathcal{N}(\log(1+A_t)),
        \mathcal{N}(|\log(1+A_t)-\log(1+A_c)|))$.
        \Statex \textit{These three inputs provide spectral evidence; they are not yet the correction map.}
        \Statex \textit{Predict a value between 0 and 1 for every channel and frequency location.}
        \State $M\gets G_\theta(Z)$.
        \Statex \textit{Use the learned gate to move suspicious amplitudes toward their clean values.}
        \State $A_{corr}\gets A_t-M\odot(A_t-A_c)$.
        \Statex \textit{Reconstruct the image using corrected amplitude and triggered phase.}
        \State $x_{corr}\gets\operatorname{IFFT}(A_{corr},P_t)$.
        \Statex \textit{Train for correct classification, small image change, sparse correction, and smoothness.}
        \State $\mathcal{L}_G\gets\operatorname{CE}(f_s(x_{corr}),y)
        +\lambda_{rec}\|x_{corr}-x\|_1
        +\lambda_{sp}\|M\|_1+\lambda_{sm}\operatorname{TV}(M)$.
        \State Update $\theta$ by minimizing $\mathcal{L}_G$.
    \EndFor
\EndFor

\State Initialize $f_r\gets f_s$ and freeze $G_\theta$.
\For{each repair epoch}
    \For{each clean minibatch $(x,y)$}
        \State Generate $t$ and $x_{corr}$ using the trained generator.
        \State Train $f_r$ on $(x,y)$, $(t,y)$, and $(x_{corr},y)$.
    \EndFor
\EndFor

\For{each non-target test image $(x,y)$}
    \State Generate $t$ and $x_{corr}$.
    \State Record predictions from $f_s$ and $f_r$ on clean, triggered,
    and corrected images.
\EndFor
\State Report clean accuracy, corrected accuracy, ASR, reconstruction error,
and correction-map visualizations.
\end{algorithmic}
\end{algorithm}

\section*{Mathematics explained in plain language}
The algorithm uses compact notation so that the complete pipeline fits in one
algorithm block. The following explanation expands the important lines.

\subsection*{1. Creating the triggered image}
The clean image is denoted by $x$. The sinusoidal trigger $S$ is a repeating
brightness pattern. The values $f_x$ and $f_y$ control how many oscillations
occur horizontally and vertically. The parameter $\alpha$ controls the size of
the added pattern:
\begin{equation}
    t=\operatorname{clip}(x+\alpha S,0,1).
\end{equation}
The function $\operatorname{clip}$ keeps every pixel in the valid image range
$[0,1]$. During poisoning, the triggered image $t_i$ is deliberately assigned
the target label $y^*$, even though its original image belongs to another
class. This teaches the suspicious classifier the hidden rule
``trigger present $\rightarrow$ target class''.

\subsection*{2. Splitting an image into amplitude and phase}
The FFT changes the representation of an image from pixels to frequencies:
\begin{equation}
    \operatorname{FFT}(x)=A_c e^{jP_c},
    \qquad
    \operatorname{FFT}(t)=A_t e^{jP_t}.
\end{equation}
Here, $A_c$ and $A_t$ describe the strength of each frequency component in the
clean and triggered images. The phase terms $P_c$ and $P_t$ describe how those
components are spatially arranged. In practical terms, amplitude says ``how
much of a frequency is present'', while phase helps determine ``where the
structure appears in the image''.

\subsection*{3. Building the generator input}
The generator does not receive only one spectrum. It receives three kinds of
evidence:
\begin{equation}
    Z=\operatorname{Concat}\left(\mathcal{N}(\log(1+A_c)),
    \mathcal{N}(\log(1+A_t)),
    \mathcal{N}\left(|\log(1+A_t)-\log(1+A_c)|\right)\right).
\end{equation}
The first block shows the clean amplitude, the second shows the triggered
amplitude, and the third shows where the two amplitudes differ. The logarithm
$\log(1+A)$ compresses very large FFT values so that weak frequencies are not
completely hidden. The operator $\mathcal{N}$ normalizes the values so that
different images and channels can be processed on a comparable scale.

The third block is evidence, not the final correction map. It says where a
change occurred; it does not decide whether that change is malicious or useful.

\subsection*{4. The generator's correction gate}
The generator processes $Z$ and produces
\begin{equation}
    M=G_\theta(Z),
    \qquad 0\leq M\leq 1.
\end{equation}
$M$ is a learned gate with one value for each channel and frequency location.
It answers the question: ``How strongly should this amplitude component be
corrected?'' It is not a binary decision and it is not the same as the
amplitude difference.

The correction rule is
\begin{equation}
    A_{corr}=A_t-M\odot(A_t-A_c),
\end{equation}
where $\odot$ means element-by-element multiplication. At one frequency:
\begin{itemize}
    \item If $M=0$, no correction is made and the triggered amplitude remains.
    \item If $M=1$, the amplitude is moved completely to the clean value.
    \item If $M=0.5$, half of the difference is removed.
\end{itemize}
Therefore, the generator can correct suspicious frequencies selectively rather
than suppressing every high-frequency component.

\subsection*{5. Reconstructing the corrected image}
The corrected amplitude is combined with the triggered phase and passed through
the inverse FFT:
\begin{equation}
    x_{corr}=\operatorname{clip}\left(\operatorname{Re}\left[
    \operatorname{IFFT}\left(A_{corr}e^{jP_t}\right)\right],0,1\right).
\end{equation}
The real part is used because numerical FFT calculations can contain tiny
imaginary rounding errors. The result is clipped again to remain a valid image.
The design changes amplitude while preserving the selected phase so that the
image structure is disturbed as little as possible.

\subsection*{6. Why the generator has four loss terms}
The generator is trained using
\begin{equation}
    \mathcal{L}_G=\mathcal{L}_{cls}
    +\lambda_{rec}\mathcal{L}_{rec}
    +\lambda_{sp}\mathcal{L}_{sp}
    +\lambda_{sm}\mathcal{L}_{sm}.
\end{equation}
Each term gives the generator a different requirement:
\begin{itemize}
    \item $\mathcal{L}_{cls}$ penalizes an incorrect prediction. The corrected
    image should receive its original clean label from the suspicious model.
    \item $\mathcal{L}_{rec}=\|x_{corr}-x\|_1$ penalizes unnecessary image
    changes and encourages the corrected image to remain close to the clean one.
    \item $\mathcal{L}_{sp}=\|M\|_1$ encourages a smaller correction gate, so
    the generator does not modify the whole spectrum without reason.
    \item $\mathcal{L}_{sm}=\operatorname{TV}(M)$ discourages unstable isolated
    values and encourages a more coherent correction map.
\end{itemize}
The coefficients $\lambda_{rec}$, $\lambda_{sp}$, and $\lambda_{sm}$ control
the trade-off between classification recovery, visual preservation, sparsity,
and map smoothness.

\subsection*{7. What the repair stage does}
After the generator is trained, its parameters are frozen. A copy of the
suspicious classifier is then fine-tuned using three versions of each training
image:
\begin{equation}
    (x,y),\qquad (t,y),\qquad (x_{corr},y).
\end{equation}
The triggered image is now given its original label $y$ rather than the
attacker's target label $y^*$. This teaches the repaired classifier that the
trigger must not change the class. The clean images preserve normal
classification ability, while the triggered and corrected images weaken the
learned backdoor association.

\subsection*{8. Reading the final ASR formula}
The attack success rate is
\begin{equation}
    \operatorname{ASR}(f)=
    \frac{\sum_i\mathbf{1}\left[\arg\max f(t_i)=y^*\right]}
    {N_{\mathrm{non-target}}}.
\end{equation}
For every non-target test image, we ask whether the triggered version is
classified as the attacker's target class. The indicator function is $1$ when
that happens and $0$ otherwise. A high suspicious ASR confirms that the
backdoor was successfully implanted. A low repaired ASR indicates that the
defense reduced the backdoor behaviour.

\section*{Spectral correction equations}
The generator receives normalized spectral evidence:
\begin{equation}
    Z=\operatorname{Concat}\left(\mathcal{N}(\log(1+A_c)),
    \mathcal{N}(\log(1+A_t)),
    \mathcal{N}\left(|\log(1+A_t)-\log(1+A_c)|\right)\right).
\end{equation}
It predicts a correction gate $M\in[0,1]^{C\times H\times W}$:
\begin{equation}
    M=G_\theta(Z).
\end{equation}
The gate moves the triggered amplitude toward the clean amplitude:
\begin{equation}
    A_{corr}=A_t-M\odot(A_t-A_c).
\end{equation}
The corrected image is reconstructed with the corrected amplitude and triggered
phase:
\begin{equation}
    x_{corr}=\operatorname{clip}\left(\operatorname{Re}\left[
    \operatorname{IFFT}\left(A_{corr}e^{jP_t}\right)\right],0,1\right).
\end{equation}

The generator objective is
\begin{equation}
    \mathcal{L}_G=\mathcal{L}_{cls}
    +\lambda_{rec}\mathcal{L}_{rec}
    +\lambda_{sp}\mathcal{L}_{sp}
    +\lambda_{sm}\mathcal{L}_{sm},
\end{equation}
where classification encourages the original label, reconstruction preserves
the clean image, sparsity limits unnecessary spectral changes, and smoothness
discourages unstable isolated values in the correction map.

\section*{FIBA amplitude-trigger alternative}
For FIBA, a reference image amplitude $A_q$ is blended into a frequency mask
$W$ while retaining the clean phase:
\begin{equation}
    A_t=(1-W)\odot A_c+W\odot\left((1-\alpha)A_c+\alpha A_q\right),
    \qquad P_t=P_c.
\end{equation}
The poisoned image is reconstructed as
\begin{equation}
    t=\operatorname{clip}\left(\operatorname{Re}\left[
    \operatorname{IFFT}\left(A_te^{jP_c}\right)\right],0,1\right).
\end{equation}
The defense then uses the same generator and correction rule.

\section*{Evaluation metric}
For a target label $y^*$, attack success rate is measured only on non-target
test images:
\begin{equation}
    \operatorname{ASR}(f)=
    \frac{\sum_{i=1}^{N_{\mathrm{non-target}}}
    \mathbf{1}\left[\arg\max f(t_i)=y^*\right]}
    {N_{\mathrm{non-target}}}.
\end{equation}

\end{document}
